%% This document is created by 
%%  Dr. Putu Harry Gunawan
%% Edited by
%%  Wisnu Adi Nurcahyo

\documentclass[a4paper,12pt,oneside]{book}
\usepackage[utf8]{inputenc}
\usepackage{sectsty}
\usepackage{graphicx}
\usepackage{epstopdf}
\usepackage{algorithm}
\usepackage{algpseudocode}
\usepackage{array}
\usepackage[table]{xcolor}
\usepackage{anysize}
\usepackage{amsmath}
\usepackage{amssymb}
\usepackage[bahasa]{babel}
\usepackage{indentfirst} %Spasi untuk paragraf pertama
\usepackage{geometry}
\usepackage{multirow}% http://ctan.org/pkg/multirow
\usepackage{hhline}% http://ctan.org/pkg/hhline
\marginsize{4cm}{3cm}{3cm}{3cm} %{left}{right}{top}{bottom}
\usepackage[compact]{titlesec} 
\usepackage{etoolbox}
\usepackage{natbib} %for Harvard syle
\usepackage[bookmarks,hypertexnames=false,debug]{hyperref}
%\usepackage[pageanchor]{hyperref}
\usepackage{bookmark}
\usepackage{dirtytalk}
% \usepackage{ccg-latex}
% \usepackage{booktabs}

% CCG
\usepackage[TABcline]{tabstackengine}

\def\tbs{\,\textbackslash}
%

\makeatletter
\patchcmd{\ttlh@hang}{\parindent\z@}{\parindent\z@\leavevmode}{}{}
\patchcmd{\ttlh@hang}{\noindent}{}{}{}

% CCG
\newcommand{\so}[1]{\ensuremath{\mbox{\it #1}^\prime}}
\newcommand\CCGFA{%
  \rlap{\kern.5\dimexpr\tabbed@gap\relax\llap{\fboxsep=0pt\colorbox{white}{$>$}}}}
\newcommand\CCGBA{%
  \rlap{\kern.5\dimexpr\tabbed@gap\relax\llap{\fboxsep=0pt\colorbox{white}{$<$}}}}
\newcommand\CCGCOOR{%
  \rlap{\kern.5\dimexpr\tabbed@gap\relax\llap{\fboxsep=0pt\colorbox{white}{\footnotesize $<\&>$}}}}
\newcommand\CCGTR{%
  \rlap{\kern.5\dimexpr\tabbed@gap\relax\llap{\fboxsep=0pt\colorbox{white}{\footnotesize $>\text{T}$}}}}
%
\makeatother

% CCG
\setlength\TABruleshift{\dimexpr.5\ht\strutbox-.5\dp\strutbox}
\setstacktabbedgap{1ex}
\TABstackTextstyle{\itshape}
\setstackgap{S}{2pt}
%

\chapterfont{\centering}
\newcommand{\bigsize}{\fontsize{16pt}{14pt}\selectfont}
\chapterfont{\centering\bigsize\bfseries}
\sectionfont{\large\bfseries}
\usepackage{tikz}
\usetikzlibrary{shapes.geometric, arrows}
%\renewcommand{\chaptertitle}{BAB}
\renewcommand{\thechapter}{\Roman{chapter}}
\renewcommand\thesection{\arabic{chapter}.\arabic{section}}
\renewcommand\thesubsection{\thesection.\arabic{subsection}}
\renewcommand{\theequation}{\arabic{chapter}.\arabic{equation}}
\renewcommand{\thefigure}{\arabic{chapter}.\arabic{figure}}
\renewcommand{\thetable}{\arabic{chapter}.\arabic{table}}

\renewcommand\bibname{Daftar Pustaka}
\addto{\captionsbahasa}{\renewcommand{\bibname}{Daftar Pustaka}}
\usepackage{fancyhdr}
\pagestyle{fancy}
\lhead{}
\chead{}
\rhead{}
\lfoot{}
\cfoot{\thepage}
\rfoot{}
\renewcommand{\headrulewidth}{0pt}

\makeatletter

%%%%%%%%%%%%%%%%%%%%%%%%%%%%%%%%%%%%%%%%%%%%%%%%%%%%%%%%%%%%
%
%  Berikut adalah data-data yang wajib diisi oleh mahasiswa
%
%%%%%%%%%%%%%%%%%%%%%%%%%%%%%%%%%%%%%%%%%%%%%%%%%%%%%%%%%%%%

% Judul dalam bahasa Indonesia
\title{
  Membangun sebuah Combinatory Categorial Grammar (CCG) Supertagger
  Berbasis Maximum Entropy untuk Bahasa Indonesia
}\let\Title\@title

% Judul dalam bahasa Inggris
\newcommand{\EngTitle}{
  Building a Combinatory Categorial Grammar (CCG) Supertagger
  Based on the Maximum Entropy for Bahasa Indonesia
}

% Nama Penulis
\author{
  Wisnu Adi Nurcahyo
}\let\Author\@author

% NIM Penulis
\newcommand{\NIM}{1301160479}

% Prodi Penulis
\newcommand{\Prodi}{Informatika}

% Kelompok Keahlian TA
\newcommand{\KK}{Intelligence, Computing, and Multimedia (ICM)}

% Gelar
\newcommand{\Gelar}{Sains Komputasi}

% Tahun TA
\date{2019}\let\Date\@date

% Kelas TA
\newcommand{\Kelas}{NLP}

% Tanggal Pengesahan
\newcommand{\Tanggal}{13}

% Bulan Pengesahan
\newcommand{\Bulan}{November}

% Nama Pembimbing Satu
\newcommand{\PembimbingSatu}{Dr. Ade Romadhony, S.T., M.T.}

% NIP Pembimbing Satu
\newcommand{\NIPPembimbingSatu}{06840042}

% Nama Pembimbing Dua
\newcommand{\PembimbingDua}{Dr. Calon Pembimbing 2, M.Kom}

% NIP Pembimbing Dua
\newcommand{\NIPPembimbingDua}{123456}

% Nama Kaprodi
\newcommand{\Kaprodi}{Niken Dwi Wahyu Cahyani, S.T., M.Kom. PhD}

% NIP Kaprodi
\newcommand{\NIPKaprodi}{xxxxxxx}


\newif\ifPembimbingHanyaSatu

%%%% WARNING kode berikut ini diaktifkan Jika pembimbing hanya satu %%%%

\PembimbingHanyaSatutrue

%%%%%%%%%%%%%%%%%%%%%%%%%%%%%%%%%%%%%%%%%%%%%%%%%%%%%%%%%%%%%%%%%%

\makeatother
\linespread{1}


\begin{document}
%%%%%%%%%%%%%%%%%%%%%%%%%%%%%%%%%%%%%%%%%%%%%%%%%%%%%%%%%%
% Dimulai dengan cover, lembar persetujuan, Abstrak, dll
%%%%%%%%%%%%%%%%%%%%%%%%%%%%%%%%%%%%%%%%%%%%%%%%%%%%%%%%%%%
\pagenumbering{roman} 
\begin{titlepage}
\thispagestyle{empty}
%\vspace*{0.7cm}
\pdfbookmark{Cover}{ }
{\centering
\large
{\bigsize\bf \Title}\\
\vspace{ 2cm}
\rm
\textbf{Proposal Tugas Akhir}\\
\vspace{0.5 cm}
\textbf{Kelas TA \Kelas}\\
\vspace{0.5 cm}
\textbf{\Author}\\ \textbf{NIM: \NIM}\\ 

\vspace{1.5 cm}

\begin{figure}[h]
{\centering {\includegraphics[scale=0.17]{Tel-U-Logo}}\par}
\end{figure}

\vspace{2 cm}
{\bigsize\textbf{Program Studi Sarjana \Prodi}\\
\vspace{0.5 cm}
\textbf{Fakultas Informatika}\\
\vspace{0.5 cm}
\textbf{Universitas Telkom}\\
\vspace{0.5 cm}
\textbf{Bandung}\\
\vspace{0.5 cm}
\textbf{\Date}\\}
}
\pagebreak
\thispagestyle{empty}
\pdfbookmark{Lembar Persetujuan}{ }
{\centering
\textbf{\large Lembar Persetujuan}\\
\vspace{0.5cm}
\textbf{\Title}\\
\vspace{0.5cm}
\textbf{\textit{\EngTitle}}\\
\vspace{0.5cm}
\textbf{\Author}\\
\textbf{NIM: \NIM}\\
\vspace{1cm}


{ Proposal ini diajukan sebagai usulan pembuatan tugas akhir pada\\ Program Studi Sarjana \Prodi\\ Fakultas Informatika Universitas Telkom}\\

\vspace{0.5cm}

{Bandung, \Tanggal\quad \Bulan \quad \Date}\\
{Menyetujui}\\

\vspace{0.5cm}
      \ifPembimbingHanyaSatu
             \begin{center}
            Calon Pembimbing 1
            \end{center}
       \else
          \begin{center}
              \begin{tabular}{  m{8cm}  m{8cm} }
              Calon Pembimbing 1 & Calon Pembimbing 2
              \end{tabular}
           \end{center}
      \fi
\begin{center}
\vspace{2cm}
\ifPembimbingHanyaSatu
     \underline{\PembimbingSatu} \\ 
      NIP: \NIPPembimbingSatu
\else
\begin{tabular}{  m{8cm}  m{8cm} }
\underline{\PembimbingSatu} & \underline{\PembimbingDua} \\ 
NIP: \NIPPembimbingSatu & NIP: \NIPPembimbingDua
\end{tabular}
\fi
\end{center}
\vspace{0.5cm}

}
\pagebreak
\end{titlepage}
\phantomsection
\addcontentsline{toc}{chapter}{Abstrak}
%\pdfbookmark{\abstractname}{Abstrak Indonesia}
\chapter*{Abstrak}

% Dalam pemrosesan bahasa alami, combinatory categorial grammar (CCG) merupakan salah satu
% formalisme tata bahasa yang dapat digunakan untuk membangun sebuah \textit{parser} yang umumnya
% dikenal sebagai CCG \textit{parser}.
% CCG \textit{parser} dapat digunakan untuk berbagai macam keperluan dalam pemrosesan bahasa alami.
% Sebagai contoh, CCG \textit{parser} dapat digunakan untuk memperoleh informasi
% (\textit{information extraction}) dari suatu kalimat yang kemudian membentuk sebuah \textit{query}.
% Agar dapat bekerja, CCG \textit{parser} membutuhkan CCG \textit{lexicon}.
% CCG \textit{lexicon} diperoleh dari proses yang bernama \textit{supertagging}.
% \textit{Supertagging} adalah proses pelabelan suatu token kata terhadap \textit{supertag}-nya.
% Perangkat lunak yang melakukan \textit{supertagging} disebut sebagai \textit{supertagger}.
% Demikian itu, \textit{supertagging} merupakan langkah pertama yang perlu dilakukan sebelum
% membangun sebuah CCG \textit{parser}.
% \textit{Supertagger} yang dibangun dalam tugas akhir ini dimaksudkan sebagai produsen CCG
% \textit{lexicon} bahasa Indonesia untuk riset-riset yang berkenaan dengan CCG di masa yang akan datang.

TBA.

\vspace{0.5 cm}
\begin{flushleft}
{\textbf{Kata Kunci:}
  natural language processing, combinatory categorial grammar, annotation tool
}
\end{flushleft}

\cleardoublepage
\phantomsection
\addcontentsline{toc}{chapter}{Daftar Isi}
%\pdfbookmark{\contentsname}{Contents}
\tableofcontents

%%%%%%%%%%%%%%%%%%%%%%%%%%%%%%%%%%%%%%%%%%%%%%%%%
% Bab bab penulisan
%%%%%%%%%%%%%%%%%%%%%%%%%%%%%%%%%%%%%%%%%%%%%%%%
\cleardoublepage
\pagenumbering{arabic}
\chapter{Pendahuluan}
\section{Latar Belakang}

Riset pemrosesan bahasa alami untuk bahasa Indonesia saat ini masih terbilang sedikit.
Bahkan, masih banyak area riset yang belum tersentuh seperti contohnya
\textit{combinatory categorial grammar} (CCG).
Sementara itu, riset mengenai CCG untuk bahasa Inggris sudah cukup matang.
Adapun untuk bahasa lainnya (seperti bahasa Vietnam) sudah mulai menggunakan CCG di dalam
penelitiannya \cite{nguyen2019vietnamese}.
Agar dapat menerapkan CCG di dalam aplikasi yang dibangun, \textit{tools} seperti
CCG \textit{parser} dan CCG \textit{supertagger} harus tersedia terlebih dahulu.
Masing-masing dari \textit{tools} tersebut memerlukan \textit{dataset} agar dapat memberikan
hasil yang akurat.

Umumnya terdapat dua cara yang paling sering digunakan untuk mengembangkan CCG \textit{supertagger}
maupun CCG \textit{parser} bahasa lokal yaitu (1) membangun \textit{dataset} CCG \textit{supertag}
secara manual maupun semi-otomatis atau (2) melakukan transfer \textit{dataset} dari CCGbank
(atau dari sumber lainnya) ke dalam bahasa lokal dengan cara melakukan alih bahasa dan bila perlu
melakukan penyesuaian untuk \textit{supertag}-nya \cite{hockenmaier-steedman-2007-ccgbank}.
Proses pembangunan \textit{dataset} umumnya menggunakan bantuan \textit{annotation tool} agar
proses anotasinya menjadi lebih mudah.
Salah satu \textit{annotation tool} yang dapat digunakan adalah
CCGweb \cite{evang-etal-2019-ccgweb}.

Tugas akhir dengan judul \say{\Title} berusaha untuk membangun alat anotasi CCG baru dengan
UI/UX yang lebih baik dari CCGweb.
Selain itu, dengan bantuan NLTK alat anotasi ini dapat melakukan \textit{generate} untuk
CCG \textit{derivation}-nya kemudian pengguna juga dapat mengubah \textit{derivation}-nya
apabila diperlukan.
Tujuan dari dibangunnya alat anotasi CCG ini adalah untuk mempermudah proses anotasi yang
repetitif.
Selanjutnya, \textit{dataset} CCG pertama untuk bahasa Indonesia diharapkan dapat dipublikasikan.
\pagebreak

\section{Perumusan Masalah}
Rumusan masalah yang akan diangkat yaitu:
\begin{enumerate}
    \item Bagaimana proses pembangunan alat anotasi untuk CCG?
    \item Apakah CCGtown dapat digunakan dengan mudah?
\end{enumerate}
\section{Tujuan}
Tujuan yang diharapkan dapat tercapai oleh tugas akhir ini yaitu:
\begin{enumerate}
    \item Membangun dan merilis alat anotasi CCG semi-otomatis.
    \item Dapat digunakan untuk membangun \textit{dataset} CCG.
\end{enumerate}
\section{Rencana Kegiatan}
Rencana kegiatan yang akan dilakukan adalah sebagai berikut:
\begin{itemize}
    \item Studi literatur
    \item Studi \textit{tools} yang tersedia
    \item Studi bahasa pemrograman yang akan digunakan
    \item Perancangan sistem \textit{annotation tool} CCGtown
    \item Membangun \textit{annotation tool} CCGtown
    \item Menguji kemudahan dalam menggunakan CCGtown
\end{itemize}

% I don't know why but this table is not showing up
\section{Jadwal Kegiatan}
Laporan proposal ini akan dijadwalkan sesuai dengan tabel \ref{ScheduleTable}.

\begin{table}
  \centering
    \caption{Jadwal kegiatan proposal tugas akhir.}
  \label{ScheduleTable}
  \begin{tabular}{|c|m{2.5cm}|m{0.01cm}|m{0.01cm}|m{0.01cm}|m{0.01cm}|m{0.01cm}|m{0.01cm}|m{0.01cm}|m{0.01cm}|m{0.01cm}|m{0.01cm}|m{0.01cm}|m{0.01cm}|m{0.01cm}|m{0.01cm}|m{0.01cm}|m{0.01cm}|m{0.01cm}|m{0.01cm}|m{0.01cm}|m{0.01cm}|m{0.01cm}|m{0.01cm}|m{0.01cm}|m{0.01cm}|}
    \hline
    \multirow{2}{*}{\textbf{No}} & \multirow{2}{*}{\textbf{Kegiatan}} & \multicolumn{24}{|c|}{\textbf{Bulan ke-}} \\
    \hhline{~~------------------------}
    {} & {} & \multicolumn{4}{|c|}{\textbf{1}} & \multicolumn{4}{|c|}{\textbf{2}} & \multicolumn{4}{|c|}{\textbf{3}} & \multicolumn{4}{|c|}{\textbf{4}} & \multicolumn{4}{|c|}{\textbf{5}} & \multicolumn{4}{|c|}{\textbf{6}}\\
    \hline
    1 & Studi Literatur & \cellcolor{blue!25} & \cellcolor{blue!25} & \cellcolor{blue!25} & \cellcolor{blue!25}& \cellcolor{blue!25} & \cellcolor{blue!25} & \cellcolor{blue!25} & \cellcolor{blue!25}& \cellcolor{blue!25} & \cellcolor{blue!25} & \cellcolor{blue!25} & \cellcolor{blue!25}& \cellcolor{blue!25} & \cellcolor{blue!25} & \cellcolor{blue!25} & \cellcolor{blue!25}& \cellcolor{blue!25} & \cellcolor{blue!25} & \cellcolor{blue!25} & \cellcolor{blue!25}& \cellcolor{blue!25} & \cellcolor{blue!25} & \cellcolor{blue!25} & \cellcolor{blue!25}\\
    \hline
    2 & Studi \textit{Tools} yang Tersedia & \cellcolor{blue!25} & \cellcolor{blue!25} & \cellcolor{blue!25} & \cellcolor{blue!25} &  \cellcolor{blue!25} &  \cellcolor{blue!25} &  \cellcolor{blue!25} &  \cellcolor{blue!25} & {} & {} & {} & {}& {} & {} & {} & {}& {} & {} & {} & {}& {} & {} & {} & {}\\
    \hline
    3 & Studi Bahasa Pemrograman dan Framework-nya & {} & {} & {} & {} & \cellcolor{blue!25} & \cellcolor{blue!25} & \cellcolor{blue!25} & \cellcolor{blue!25} & {} & {} & {} & {}& {} & {} & {} & {}& {} & {} & {} & {}& {} & {} & {} & {}\\
    % \hline
    % 4 & Pengumpulan Data & {} & {} & {} & {} & \cellcolor{blue!25} & \cellcolor{blue!25} & \cellcolor{blue!25} & \cellcolor{blue!25} & \cellcolor{blue!25} & \cellcolor{blue!25} & \cellcolor{blue!25} & \cellcolor{blue!25} & {} & {} & {} & {}& {} & {} & {} & {} & {} & {} & {} & {}\\
    \hline
    5 & Analisis dan Perancangan Sistem &  {} & {} & {} & {}  & \cellcolor{blue!25} & \cellcolor{blue!25} & \cellcolor{blue!25} & \cellcolor{blue!25} & \cellcolor{blue!25} & \cellcolor{blue!25} & \cellcolor{blue!25} & \cellcolor{blue!25} & {} & {} & {} & {}& {} & {} & {} & {}& {} & {} & {} & {}\\
    \hline
    6 & Implementasi Sistem &  {} & {} & {} & {} & {} & {} & {} & {}& \cellcolor{blue!25} & \cellcolor{blue!25} & \cellcolor{blue!25} & \cellcolor{blue!25} & \cellcolor{blue!25} & \cellcolor{blue!25} & \cellcolor{blue!25} & \cellcolor{blue!25} & {} & {} & {} & {}& {} & {} & {} & {}\\
    \hline
    7 & Analisa Hasil Implementasi &  {} & {} & {} & {} & {} & {} & {} & {}& {} & {} & {} & {} & \cellcolor{blue!25} & \cellcolor{blue!25} & \cellcolor{blue!25} & \cellcolor{blue!25} & \cellcolor{blue!25} & \cellcolor{blue!25} & \cellcolor{blue!25} & \cellcolor{blue!25} & {} & {} & {} & {}\\
    \hline
    8 & Penulisan Laporan & {} & {} & {} & {} & \cellcolor{blue!25} & \cellcolor{blue!25} & \cellcolor{blue!25} & \cellcolor{blue!25}& \cellcolor{blue!25} & \cellcolor{blue!25} & \cellcolor{blue!25} & \cellcolor{blue!25}& \cellcolor{blue!25} & \cellcolor{blue!25} & \cellcolor{blue!25} & \cellcolor{blue!25}& \cellcolor{blue!25} & \cellcolor{blue!25} & \cellcolor{blue!25} & \cellcolor{blue!25}& \cellcolor{blue!25} & \cellcolor{blue!25} & \cellcolor{blue!25} & \cellcolor{blue!25}\\
    \hline
  \end{tabular}
\end{table}

%
\chapter{Kajian Pustaka}

\section{Categorial Grammar}
Categorial Grammar (CG) merupakan sebuah istilah yang mencakup beberapa formalisme terkait yang diajukan
untuk sintaks dan semantik dari bahasa alami serta untuk bahasa logis dan matematis \cite{Steedman92catg}.
Karakteristik yang paling terlihat dari CG adalah bentuk ekstrim dari leksikalismenya di mana beban utama
(atau bahkan seluruh beban) sintaksisnya ditanggung oleh leksikon.
Konstituen tata bahasa dalam \textit{categorial grammar} dan khususnya semua leksikal diasosiasikan
dengan suatu \textit{type} atau \say{\textit{category}} (dalam \textit{category theory}) yang
mendefinisikan potensi mereka untuk dikombinasikan dengan konstituen lain untuk menghasilkan konstituen
majemuk.
\textit{Category} tersebut adalah salah satu dari sejumlah kecil \textit{category} dasar (seperti NP)
atau \textit{functor} (dalam \textit{category theory}).
Dalam hal ini, \textit{category} dapat diartikan sebagai \textit{syntactic type} dari suatu kata.

Secara formal, \textit{syntactic type} didefinisikan sebagai himpunan bagian dari suatu
\textit{semigroup} $M$ yang tunduk pada tiga operasi yaitu \ref{catg:syn:1},
\ref{catg:syn:2}, dan \ref{catg:syn:3} dimana $A$, $B$, dan $C$ merupakan himpunan bagian dari $M$
\cite{Lambek1988}. Adapun $A \cdot B$ dibaca $A$ \textit{times} $B$, $C/B$ dibaca $C$ \textit{over}
$B$, dan $A\backslash{}C$ dibaca $A$ \textit{under} $C$. Selanjutnya, dapat dilihat bahwasannya
untuk semua $A, B, C \subseteq M$ sehingga kita dapatkan \ref{catg:syn:4} dan \ref{catg:syn:5}.
Terakhir, persamaan \ref{catg:syn:6} dapat diabaikan apabila dihadapkan dengan
\textit{multiplicative system} yang tidak asosiatif. Sementara itu, apabila \textit{semigroup}-nya
merupakan sebuah \textit{monoid} dengan identitas $1$ maka kita dapatkan \ref{catg:syn:7} dimana
$I = \{1\}$.

\begin{align}
  \begin{split}\label{catg:syn:1}
    A \cdot B & = \{x \cdot y \in M \mid x \in A \land y \in B\}
  \end{split}\\
  \begin{split}\label{catg:syn:2}
    C/B & = \{x \in M \mid \forall_{y \in B} x \cdot y \in C\}
  \end{split}\\
  \begin{split}\label{catg:syn:3}
    A\backslash{}C & = \{y \in M \mid \forall_{x \in A} x \cdot y \in C\}
  \end{split}
\end{align}

\begin{align}
  \begin{split}\label{catg:syn:4}
    A \cdot B \subseteq C & \;\;\;\;\text{jika dan hanya jika}\;\;\;\; A \subseteq C/B
  \end{split}\\
  \begin{split}\label{catg:syn:5}
    A \cdot B \subseteq C & \;\;\;\;\text{jika dan hanya jika}\;\;\;\; B \subseteq A\backslash{}C
  \end{split}
\end{align}

\begin{align}
  \begin{split}\label{catg:syn:6}
    (A \cdot B) \cdot C = A \cdot (B \cdot C)
  \end{split}\\
  \begin{split}\label{catg:syn:7}
    I \cdot A = A = A \cdot I
  \end{split}
\end{align}

Ada beberapa notasi berbeda untuk \textit{category} dalam merepresentasikan \textit{directional}-nya.
Notasi yang paling umum digunakan adalah \say{\textit{slash notation}} yang dipelopori oleh Bar-Hilel,
Lambek, dan kemudian dimodifikasi dalam kelompok teori yang dibedakan sebagai tata bahasa
\say{\textit{combinatory}} \textit{categorial grammar} (CCG).
Sebagai contoh, \textit{category} $\text{(S$\backslash$NP)/NP}$ merupakan suatu \textit{functor} yang
memiliki dua buah notasi \textit{slash} yaitu $\backslash$ dan $/$.
Masing-masing notasi \textit{slash} tersebut merepresentasikan \textit{directionality} yang berbeda.
Notasi \textit{forward slash}, $/$, mengindikasikan bahwa argumen dari suatu \textit{functor}
$\text{X}/\text{Y}$ ada di bagian kanan atau dengan kata lain $\text{Y}$.
Adapun \textit{backward slash}, $\backslash$, mengindikasikan bahwa argumen dari suatu \textit{functor}
$\text{X}\backslash\text{Y}$ ada di bagian kiri atau dengan kata lain $\text{X}$.
Demikian itu, penggunaan notasi \textit{slash} yang tepat sangat penting dikarenakan hal ini dapat
mempengaruhi konstituen dari hasil \say{kombinasi} \textit{category}-nya.


\section{Combinatory Categorial Grammar}\label{kajian-ccg}
Combinatory Categorial Grammar (CCG) merupakan salah satu formalisme tata bahasa yang gaya aturannya
diturunkan dari \textit{categorial grammar} dengan beberapa penambahan aturan dan istilah baru
\cite{Steedman96avery}.
Di CCG, \textit{category} dapat dipasangkan dengan \textit{semantic representation}.
Dalam hal ini, \textit{semantic representation} yang dimaksud adalah abstraksi fungsi lambda
(dalam \textit{lambda calculus}, \textit{lambda function}).
Sebagai contoh, \textit{category} $\text{(S$\backslash$NP)/NP}$ dapat dipasangkan dengan fungsi lambda
$\lambda{x. fx}$ sehingga dapat ditulis menjadi $\text{(S$\backslash$NP)/NP} : \lambda{x. fx}$.
Adapun pemetaan dari suatu token kata ke \textit{category}-nya menggunakan notasi $\vdash$.
Sebagai contoh, anggap saja kita memiliki kamus pemetaan seperti pada Gambar \ref{ccg:mapping:1}.
Apabila kita memiliki kalimat \say{Pamungkas dan Setyo menyukai rendang}, maka kita dapatkan:

\begin{figure}\centering\small
  \begin{align*}
    \text{Pamungkas} &\ \vdash\ \text{NP}:\ \so{pamungkas}\\
    \text{Setyo} &\ \vdash\ \text{NP}:\ \so{setyo}\\
    \text{dan} &\ \vdash\ \text{CONJ}:\ \lambda x.\lambda y.\lambda f.\ (f\ x) \land (f\ y)\\
    \text{menyukai} &\ \vdash\ \text{(S{$\backslash$}NP)/NP}:\ \lambda x.\lambda y.\ suka(y, x)\\
    \text{rendang} &\ \vdash\ \text{NP}:\ \so{rendang}
  \end{align*}
  \caption{Kamus yang memetakan token kata ke bentuk CCG \textit{lexicon}-nya.}
  \label{ccg:mapping:1}
\end{figure}

\begin{center}
  \bgroup
  \catcode`!=\active \def!{\upshape}
  \catcode`?=\active \def?#1{\makebox[0pt]{#1}}
  \catcode`^=\active \def^#1{\footnotesize{#1}}
  \catcode`*=\active \def*#1{\scriptsize{#1}}
  \tabbedShortstack{
    !^Pamungkas & & !^dan & & !^Setyo & & !^menyukai & & !^rendang &\\
    \TABcline{1,3,5,7,9}
    !^{$\text{NP}$} & &
      !^{$\text{CONJ}$} & &
      !^{$\text{NP}$} & &
      !^{$\text{(S$\backslash$NP)/NP}$} & &
      !^{$\text{NP}$} &\\
    !{*: \so{pamungkas}} & &
      !{*: $\lambda x.\lambda y.\lambda f.\ (f\ x) \land (f\ y)$} & &
      !{*: \so{setyo}} & &
      !{*: $\lambda x.\lambda y.\ suka(y, x)$} & &
      !{*: \so{rendang}} &
  }
  \egroup
\end{center}

Ada beberapa operasi yang dapat dilakukan dalam CCG. \textit{Operand} dari operasi
yang dimaksud adalah \textit{category}. Berdasarkan contoh di atas, akan ada tiga
operasi yang dijalankan yaitu \textit{coordination}, \textit{forward application},
dan \textit{type rising}.
Untuk mendapatkan hasil yang diinginkan, kita lakukan \textit{type rising} sebelum
\textit{forward application} di akhir.
Sehingga, kita dapatkan:

\begin{center}
  \bgroup
  \catcode`!=\active \def!{\upshape}
  \catcode`?=\active \def?#1{\makebox[0pt]{#1}}
  \catcode`^=\active \def^#1{\footnotesize{#1}}
  \catcode`*=\active \def*#1{\scriptsize{#1}}
  \tabbedLongunderstack{
    !^Pamungkas & & !^dan & & !^Setyo & & !^menyukai & & !^rendang &\\
    \TABcline{1,3,5,7,9}
    !^{$\text{NP}$} & &
      !^{$\text{CONJ}$} & &
      !^{$\text{NP}$} & &
      !^{$\text{(S$\backslash$NP)/NP}$} & &
      !^{$\text{NP}$} &\\
    !{*: \so{pamungkas}} & &
      !{*: $\lambda x.\lambda y.\lambda f.\ (f\ x) \land (f\ y)$} & &
      !{*: \so{setyo}} & &
      !{*: $\lambda x.\lambda y.\ suka(y, x)$} & &
      !{*: \so{rendang}} &\\
    \TABrule & \TABrule &
      \TABrule & \TABrule &
      \TABrule\CCGCOOR & &
      \TABrule & \TABrule &
      \TABrule\CCGFA &\\
    & &
      !?{^{$\text{NP}$}}
      \ \ \ \ \ \ \ \ \ 
      & & & &
      \ \ \ \ \ \ \ 
      !?{^{$\text{S$\backslash$NP}$}}
      & & &\\
    & &
      ?{*: $\lambda f.\ (f\ \so{pamungkas}) \land (f\ \so{setyo})$}
      \ \ \ \ \ \ \ \ \ 
      & & & &
      \ \ \ \ \ \ \ 
      ?{*: $\lambda y.\ suka(y, \so{rendang})$}
      & & &\\
    \TABrule & \TABrule &
      \TABrule & \TABrule &
      \TABrule\CCGTR & &
      & & &\\
    & &
      !?{^{$\text{S/(S$\backslash$NP)}$}}
      \ \ \ \ \ \ \ \ \ 
      & & & & & & &\\
    & &
      ?{*: $\lambda f.\ (f\ \so{pamungkas}) \land (f\ \so{setyo})$}
      \ \ \ \ \ \ \ \ \ 
      & & & & & & &\\
    \TABrule & \TABrule &
      \TABrule & \TABrule &
      \TABrule & \TABrule &
      \TABrule & \TABrule &
      \TABrule\CCGFA &\\
    & & &
      !?{^{$\text{S}$}}
      & & & & & &\\
    & & &
      ?{*: $suka(\so{pamungkas}, \so{rendang}) \land suka(\so{setyo}, \so{rendang})$}
      & & & & & &
  }
  \egroup
\end{center}

Berdasarkan hasil evaluasi tersebut, kita dapatkan \textit{query} \ref{ccg:query:1}
yang diperoleh dari kalimat \say{Pamungkas dan Setyo menyukai rendang}.
Demikian itu, komputer dapat melakukan komputasi berdasarkan \textit{query} yang telah diperoleh.
Kegiatan tersebut merupakan apa yang disebut dengan CCG \textit{parsing}.
Untuk dapat melakukan parsing, CCG \textit{lexicon} diperlukan.
Untuk mendapatkan CCG \textit{lexicon} kita dapat menggunakan CCG \textit{supertagger}
yang akan melakukan pelabelan suatu token kata ke CCG \textit{lexicon} berdasarkan
pemetaannya.

\begin{equation}\label{ccg:query:1}
  suka(\so{pamungkas}, \so{rendang}) \land suka(\so{setyo}, \so{rendang})
\end{equation}


\section{Lambda Calculus}
\textit{Lambda calculus} ({$\lambda$}\textit{-calculus}) merupakan sebuah formalisme yang dikembangkan
oleh Alonzo Church sebagai alat yang digunakan untuk memahami konsep komputasi yang efektif
\cite{DBLP:journals/corr/Rojas15}.
Formalisme {$\lambda$}\textit{-calculus} cukup populer dan bahkan dijadikan sebagai pondasi teori bagi
paradigma pemrograman \textit{functional programming}.
Konsep utama dari {$\lambda$}\textit{-calculus} adalah apa yang disebut dengan \textit{expression}.
Suatu \textit{expression} dalam {$\lambda$}\textit{-calculus} terdiri dari tiga bagian yaitu
\textit{lambda notation} ({$\lambda$}), \textit{argument} (seperti $a$, $b$, $c$, $x$, dan lain-lain),
dan \textit{body} yang dipisahkan dengan tanda titik.
Sebagai contoh, fungsi lambda ${\lambda}x. x$ merupakan sebuah fungsi identitas yang mengambil
argumen $x$ kemudian mengembalikan nilai $x$ itu sendiri.
Dalam hal ini, terlihat bahwa notasi {$\lambda$} merupakan sebuah penanda bagi suatu fungsi lambda.
Kemudian, pengubah $x$ setelah notasi {$\lambda$} merupakan argumen dari fungsi tersebut.
Selanjutnya, tanda titik merupakan pemisah antara \textit{head} dan \textit{body} fungsi lambda.
Terakhir, setelah tanda titik adalah \textit{body} dari suatu fungsi lambda yang mana berupa
\textit{expression}.

Untuk mempermudah pemahaman, {$\lambda$}\textit{-calculus} dapat diperlakukan seperti fungsi tanpa
nama. Sebagai contoh, fungsi lambda $({\lambda}x. x + 5)$ apabila diberikan nilai $2$ sehingga
menjadi $({\lambda}x. x + 5) 2$ akan dievaluasi menjadi ${\lambda}(2). (2) + 5$.
Demikian itu, nilai yang dikembalikan oleh fungsi tersebut adalah $7$.
Sama seperti fungsi pada umumnya, konsep ini bernama \textit{substition} (substitusi).
Memahami {$\lambda$}\textit{-calculus} dirasa perlu berhubung dalam tugas akhir ini
{$\lambda$}\textit{-calculus} digunakan sebagai bentuk formal di \textit{category}
dalam konteks CCG \textit{lexicon}. Meskipun {$\lambda$}\textit{-calculus} tidak sesederhana
yang dijelaskan sebelumnya, setidaknya memahami {$\lambda$}\textit{-calculus} seperti ini
sudah cukup untuk dapat membangun \textit{supertagger} yang ada di tugas akhir ini.


\section{CCGweb}
CCGweb\footnote{\url{https://github.com/texttheater/ccgweb}} merupakan
\textit{open source graphical annotation tool} pertama untuk CCG \cite{evang-etal-2019-ccgweb}.
Aplikasinya berbasis web dan dibangun dengan menggunakan bahasa pemrograman
Python, PHP, dan JavaScript.
Fitur yang paling menarik dari \textit{graphical annotation tool} adalah What You See Is What
You Get (WYSIWYG) yang mana berupa kemampuan untuk me-\textit{render} CCG \textit{derivation}
sesuai dengan apa yang kita lihat.
Maksudnya, CCG \textit{derivation} akan ditampilkan horizontal sesuai dengan panjang kalimatnya
kemudian hasil \textit{derivation}-nya ditampilkan vertikal seperti contoh pada
Bagian \ref{kajian-ccg}.

Untuk dapat menggunakan CCGweb, kita perlu melakukan instalasi terlebih dahulu.
Selanjutnya barulah kita dapat menambahkan kalimat-kalimat yang ingin dianotasi.
Satu instalasi CCGweb hanya dapat digunakan untuk satu proyek anotasi sehingga
apabila kita memiliki lebih dari satu proyek maka kita perlu melakukan instalasi
CCGweb yang baru.
Demikian itu, CCGtown\footnote{\url{https://github.com/wisn/ccgtown}} hadir dengan
fitur \textit{multi-project} dan tanpa perlu melakukan instalasi di komputer lokal
karena aplikasinya \textit{hosted} sehingga dapat diakses kapan pun.

%
\chapter{Perancangan Sistem}

% Sebelum dapat membangun CCG \textit{supertagger}, sebuah \textit{dataset} diperlukan untuk melatih model
% \textit{classifier} MaxEnt.
% Berhubung \textit{dataset} CCG \textit{supertag} untuk bahasa Indonesia belum tersedia, tidak ada pilihan
% lain selain membuat \textit{dataset}-nya terlebih dahulu.
% Salah satu cara untuk membuat \textit{dataset} adalah dengan memanfaatkan POS \textit{tag} dari suatu
% lema yang kemudian diberikan \textit{supertag}-nya secara manual.
% Setelah \textit{dataset} memiliki \textit{supertag} yang cukup, barulah proses \textit{train} MaxEnt
% dapat dilakukan.
% Model yang telah dilatih akan disimpan ke dalam format JSON khusus agar ke depannya dapat digunakan
% tanpa perlu melakukan \textit{train} lagi.
% Setelah itu, diperlukan program yang dapat memanfaatkan model yang telah dilatih untuk membuktikan
% bahwasannya model yang telah dibangun dapat menghasilkan \textit{supertag} yang sesuai.
% Adapun langkah terakhirnya adalah dengan menyediakan dukungan \textit{web application} berbasis
% RESTful API agar siapapun dapat memanfaatkan \textit{supertagger} ini tanpa perlu terhalang oleh
% batasan bahasa pemrograman yang digunakan.

% \section{Pembuatan Dataset}

% \textit{Dataset} untuk melakukan \textit{train} CCG sayangnya belum tersedia.
% Karenanya, kita akan membuatkan \textit{dataset} secara semi otomatis dengan memanfaatkan POS
% \textit{tagger} untuk bahasa Indonesia.
% Secara formal, \textit{flowchart} untuk proses pembuatan \textit{dataset} dapat dilihat pada
% Gambar \ref{flowchart:dataset:1}.

% \begin{figure}[h!]\centering
% 	\begin{tikzpicture}[node distance=2cm]
% 		\node (start) [cloud] {Start};
% 		\node (in) [io, below of=start, yshift=0.5cm] {\textit{Input} teks bahasa Indonesia};
% 		\node (pro1) [process, below of=in, yshift=0.5cm] {POS \textit{tagging}};
% 		\node (pro2) [process, below of=pro1, yshift=0.5cm] {Jadikan \textit{queue}};
% 		\node (dec) [decision, below of=pro2, yshift=-1cm] {\textit{Queue} kosong?};
% 		\node (pro3) [process, below of=dec, yshift=-1cm] {Transformasi POS \textit{tag}};
% 		\node (out) [io, below of=pro3, yshift=0.5cm] {\textit{Output} CCG \textit{supertag}};
% 		\node (stop) [cloud, below of=out, yshift=0.5cm] {Stop};
% 		\draw [arrow] (start) -- (in);
% 		\draw [arrow] (in) -- (pro1);
% 		\draw [arrow] (pro1) -- (pro2);
% 		\draw [arrow] (pro2) -- (dec);
% 		\draw [arrow] (dec) -- node[anchor=east]{Tidak} (pro3);
% 		\draw [arrow] (pro3) -- +(-3,0) |- (dec);
% 		\draw [arrow] (dec) -- node[anchor=south]{Ya} +(4,0) |- (out);
% 		\draw [arrow] (out) -- (stop);
% 	\end{tikzpicture}
% 	\caption{Alur kerja pembuatan \textit{dataset} untuk CCG \textit{supertag}.}
% 	\label{flowchart:dataset:1}
% \end{figure}

% \subsection{Input Teks Bahasa Indonesia}

% Pada bagian \textit{input} dalam pembuatan \textit{dataset}, kita dapat mengambil teks bahasa Indonesia
% dari Indonesian Treebank\footnotemark[1]\ dan/atau dari beberapa contoh artikel yang terdapat di
% \textit{website} Wikipedia Indonesia\footnotemark[2].
% Indonesian Treebank terdapat setidaknya 1000 kalimat yang dirasa cukup untuk melakukan
% \textit{training} menggunakan model MaxEnt.
% Apa yang perlu dilakukan setelah mengambil \textit{input} adalah membersihkan bentuk \textit{tree}-nya
% untuk kemudian diambilkan kalimat yang sebenarnya.
% Caranya cukup sederhana yaitu dengan mengambil \textit{leaf} dari \textit{tree} tersebut kemudian
% disatukan di suatu pengubah dengan tipe data \textit{string}.

% Sebagai contoh, salah satu \textit{tree} yang terdapat di Indonesian Treebank dapat dilihat pada
% Gambar \ref{treebank:tree:1}.
% Apa yang dimaksud dengan \textit{leaf} pada \textit{tree} tersebut adalah $(Kera)$, $(untuk)$,
% $(*)$, $(amankan)$, $(pesta\ olahraga)$.
% Terkhusus untuk \textit{leaf} dengan bentuk spesial, seperti $(*)$, kita hilangkan sehingga teks
% yang diperoleh dari \textit{tree} tersebut adalah \say{Kera untuk amankan pesta olahraga}.
% Selanjutnya, contoh teks yang akan digunakan agar konsisten yaitu
% \say{Pamungkas dan Setyo menyukai rendang}.


% \begin{figure}\centering\small
% 	\begin{align*}
% 		&(NP\\
% 		&\ \ (NN\ (Kera))\\
% 		&\ \ (SBAR\\
% 		&\ \ \ \ (SC\ (untuk))\\
% 		&\ \ \ \ (S\ (NP-SBJ\ (*))\\
% 		&\ \ \ \ \ \ (VP\\
% 		&\ \ \ \ \ \ \ \ (VB\ (amankan))\\
% 		&\ \ \ \ \ \ \ \ (NP\ (NN\ (pesta\ olahraga)))))))
% 	\end{align*}
% 	\caption{Salah satu contoh \textit{tree} dalam Indonesian Treebank.}
% 	\label{treebank:tree:1}
% \end{figure}

% \footnotetext[1]{github.com/famrashel/idn-treebank/blob/master/Indonesian\_Treebank.bracket}
% \footnotetext[2]{id.wikipedia.org}

% \subsection{POS Tagging}

% Berdasarkan teks bahasa Indonesia yang telah diambil, kita manfaatkan \textit{tool} POS \textit{tagger}
% bahasa Indonesia untuk mendapatkan \textit{lexical category} atomik untuk masing-masing token.
% Dengan memanfaatkan POS \textit{tag} kita dapat membuat \textit{dataset} untuk CCG \textit{supertag}
% lebih mudah dibandingkan dengan memberikan \textit{tag} CCG secara manual.
% Ide dasarnya yaitu kita akan mentransformasikan POS \textit{tag} yang didapatkan menjadi CCG
% \textit{supertag} berdasarkan aturan-aturan khusus yang telah ditentukan.
% Sebagai contoh, kita dapat membuat aturan seperti mentransformasikan $\text{VB}$ menjadi
% $\text{(S$\backslash$NP)/NP}$.
% Dengan menggunakan contoh kalimat yang sama seperti di bagian sebelumnya,
% yaitu \say{Pamungkas dan Setyo menyukai rendang}, setelah menggunakan POS \textit{tagger} bahasa
% Indonesia kita dapatkan hasil sesuai dengan Gambar \ref{postag:1}.

% \begin{figure}\centering
%   \bgroup
%   \catcode`!=\active \def!{\upshape}
%   \tabbedShortstack{
%     !Pamungkas & & !dan & & !Setyo & & !menyukai & & !rendang &\\
%     \TABcline{1,3,5,7,9}
%     !{$\text{NNP}$} & &
%       !{$\text{CC}$} & &
%       !{$\text{NNP}$} & &
%       !{$\text{VB}$} & &
%       !{$\text{X}$} &
%   }
% 	\egroup
% 	\caption{Kalimat contoh dengan POS \textit{tag}-nya.}
% 	\label{postag:1}
% \end{figure}

% \subsection{Transformasi POS Tag}

% Pada bagian proses transformasi POS \textit{tag} ke bentuk CCG \textit{supertag}-nya,
% kita buatkan aturan-aturan transformasinya.
% Aturan transformasi tersebut merupakan pemetaan berbasis aturan.
% Sebagai contoh, kata \say{menyukai} memiliki POS \textit{tag} $\text{VB}$.
% Anggap saja dalam aturan transformasi terdapat pemetaan $\text{VB} \vdash \text{(S$\backslash$NP)/NP}$.
% Sehingga, kita dapatkan CCG \textit{supertag} untuk \say{menyukai} yaitu $\text{(S$\backslash$NP)/NP}$.
% Kendati demikian, masih ada bagian yang belum kita dapatkan yaitu \textit{semantic representation}-nya.
% Kita dapat menggunakan \textit{stemmer} bahasa Indonesia agar mendapatkan \textit{root words}
% dari kata \say{menyukai} yaitu \say{suka}.
% Langkah terakhirnya adalah membuatkan \textit{semantic representation}-nya berdasarkan
% \textit{root words} yang telah diperoleh sehingga kita dapatkan fungsi lambdanya yaitu
% $\lambda x.\lambda y.\ \text{suka}(y, x)$.

% Selain menggunakan \textit{stemmer}, kita dapat menggunakan \textit{morphological analyzer}.
% \textit{Morphological analyzer} salah satu kegunaannya yaitu dapat menghasilkan \textit{root words}
% sehingga dapat kita manfaatkan untuk membuat fungsi lambda.
% Untuk bahasa Indonesia, kita dapat menggunakan \textit{tool} bernama MorphInd\footnotemark[3].
% Selain MorphInd, alternatif \textit{tool} yang dapat digunakan adalah IndMA
% (Indonesian morphological analyzer).
% Namun, MorphInd dipilih karena dapat memberikan keluaran berupa morfem tersegmentasi
% \cite{larasati2011indonesian}.
% Kita dapat menggunakan MorphInd sebagai pelengkap POS \textit{tagger} untuk bahasa Indonesia.
% Hal ini agar \textit{dataset} yang dibuatkan secara semi-otomatis ini dapat memiliki kualitas yang
% baik.


% \footnotetext[3]{septinalarasati.com/morphind}

% \subsection{Output CCG Supertag}

% Keluaran dari bagian pembuatan \textit{dataset} ini adalah sebuah berkas JSON
% (JavaScript Object Notation) berisi CCG \textit{supertag} lengkap dan beberapa berkas JSON dari CCG
% \textit{supertag} yang lemanya dikelompokkan berasarkan alfabet.
% Adapun format JSON dari \textit{dataset} yang disimpan dalam berkas tersebut dapat dilihat pada Gambar
% \ref{dataset:format}.
% Untuk setiap kalimat dalam berkas tersebut direpresentasikan oleh dua objek yaitu
% (1) \say{tokens} berupa daftar token, dan
% (2) \say{supertags} berupa daftar \textit{supertag} untuk token ke-$(i, j)$ dimana $1 \leq i \leq n$
% dan $1 \leq j \leq m$ dalam sistem \textit{$1$-indexed array}.
% Sebagai contoh, kalimat \say{Pamungkas menyukai rendang} representasinya dapat dilihat pada Gambar
% \ref{dataset:example}.
% Demikian itu, kita dapat mengambil \textit{dataset} secara lengkap dengan CCG \textit{supertag}-nya.

% \begin{figure}\centering
% 	\bgroup
% 	\catcode`!=\active \def!{\upshape}
% 	\tabbedShortstack[l]{
% 		!{[} & & &\\
% 			& !{\{} & &\\
% 			& & !{\say{tokens}:} &
% 				!{[\say{$token_{1,1}$}, \say{$token_{1,2}$}, $\dots$, \say{$token_{1,m}$}],}\\
% 			& & !{\say{supertags}:} &
% 				!{[\say{$supertag_{1,1}$}, \say{$supertag_{1,2}$}, $\dots$, \say{$supertag_{1,m}$}]}\\
% 			& !{\}}, & &\\
% 			& !{\{} & &\\
% 			& & !{\say{tokens}:} &
% 				!{[\say{$token_{2,1}$}, \say{$token_{2,2}$}, $\dots$, \say{$token_{2,m}$}],}\\
% 			& & !{\say{supertags}:} &
% 				!{[\say{$supertag_{2,1}$}, \say{$supertag_{2,2}$}, $\dots$, \say{$supertag_{2,m}$}]}\\
% 			& !{\}}, & &\\
% 			& & &\\
% 			& & &\\
% 			& $\vdots$ & &\\
% 			& & &\\
% 			& & &\\
% 			& !{\{} & &\\
% 			& & !{\say{tokens}:} &
% 				!{[\say{$token_{n,1}$}, \say{$token_{n,2}$}, $\dots$, \say{$token_{n,m}$}],}\\
% 			& & !{\say{supertags}:} &
% 				!{[\say{$supertag_{n,1}$}, \say{$supertag_{n,2}$}, $\dots$, \say{$supertag_{n,m}$}]}\\
% 			& !{\}} & &\\
% 		!{]} & & &
% 	}
% 	\egroup
% 	\caption{Format JSON \textit{dataset} yang disimpan di dalam berkas.}
% 	\label{dataset:format}
% \end{figure}

% % \begin{figure}\centering
% % 	\bgroup
% % 	\catcode`!=\active \def!{\upshape}
% % 	\catcode`?=\active \def?#1{\makebox[0pt]{#1}}
% % 	\tabbedShortstack[l]{
% % 		!{[} & & &\\
% % 			& !{\{} & &\\
% % 			& & !{\say{tokens}: [} &\\
% % 			& & !{\ \ \ \ \say{Pamungkas},} &\\
% % 			& & !{\ \ \ \ \say{dan},} &\\
% % 			& & !{\ \ \ \ \say{Setyo},} &\\
% % 			& & !{\ \ \ \ \say{menyukai},} &\\
% % 			& & !{\ \ \ \ \say{rendang},} &\\
% % 			& & !{\ \ \ \ \say{.}} &\\
% % 			& & !{],} &\\
% % 			& & !{\say{supertags}: [} &\\
% % 			& & !{\ \ \ \ \say{$\text{NP: \so{pamungkas}}$},} &\\
% % 			& & !{\ \ \ \ \say{$\text{CONJ: $\lambda x.\lambda y.\lambda f.\ (f\ x) \land (f\ y)$}$},} &\\
% % 			& & !{\ \ \ \ \say{$\text{NP: \so{setyo}}$},} &\\
% % 			& & !{\ \ \ \ \say{$\text{(S$\backslash$NP)/NP: $\lambda x.\lambda y.\ suka(y, x)$}$},} &\\
% % 			& & !{\ \ \ \ \say{$\text{NP: \so{rendang}}$},} &\\
% % 			& & !{\ \ \ \ \say{$\text{Z}$}} &\\
% % 			& & !{]} &\\
% % 			& !{\}} & &\\
% % 		!{]} & & &
% % 	}
% % 	\egroup
% % 	\caption{Contoh isi dari berkas \textit{dataset} dalam format JSON.}
% % 	\label{dataset:example}
% % \end{figure}

% \begin{figure}\centering\scriptsize
% 	\bgroup
% 	\catcode`!=\active \def!{\upshape}
% 	\tabbedShortstack[l]{
% 		!{[} & & &\\
% 			& !{\{} & &\\
% 			& & !{\say{tokens}:} &
% 				!{[\say{Pamungkas}, \say{menyukai}, \say{rendang}, \say{.}],}\\
% 			& & !{\say{supertags}:} &
% 				!{[
% 					\say{$\text{NP: \so{pamungkas}}$},
% 					\say{$\text{(S$\backslash$NP)/NP: $\lambda x.\lambda y.\ suka(y, x)$}$},
% 					\say{$\text{NP: \so{rendang}}$},
% 					\say{$\text{Z}$}
% 				]}\\
% 			& !{\}} & &\\
% 		!{]} & & &
% 	}
% 	\egroup
% 	\caption{Contoh isi dari berkas \textit{dataset} dalam format JSON.}
% 	\label{dataset:example}
% \end{figure}

% \section{Mempersiapkan CCG Supertagger}

% \subsection{Melatih Model MaxEnt}

% Model yang akan digunakan oleh \textit{supertagger} ini adalah Maximum Entropy (MaxEnt).
% Model ini dipilih karena ketersediaan \textit{dataset} bahasa Indonesia yang masih sangat kurang.
% Selain itu, model MaxEnt digunakan di sebuah riset yang dilakukan oleh Stephen Clark dalam
% pembuatan \textit{supertagger}-nya.
% Bahkan, performansi \textit{supertagger} yang dikembangkan sangat baik.
% Riset tersebut pada intinya membuktikan bahwasannya MaxEnt dapat digunakan di \textit{supertagger}
% juga meskipun pada peruntukannya MaxEnt dibuat untuk POS \textit{tagger}.
% Tentunya terdapat beberapa penyesuaian yang harus dilakukan.
% Salah satunya adalah formula probabiltas yang digunakan.
% Formula yang telah disesuaikan tersebut dapat dilihat di persamaan \ref{maxent:equation:2}.

% Model yang telah dilatih akan disimpan di dalam sebuah berkas sebagai \textit{snapshot} agar apabila
% proses \textit{train} mengalami kegagalan kita tidak perlu mengulangi proses \textit{train} dari
% awal lagi.
% Dalam hal ini, untuk setiap $10$ \textit{dataset} yang telah diproses, proses penyimpanan
% \textit{snapshot} akan dilakukan.
% Adapun format \textit{snapshot} menggunakan JSON yang menyimpan informasi berupa jumlah kalimat
% terproses dan sebuah tabel yang menyimpan kalkulasi pada saat \textit{snapshot} disimpan.
% Setelah proses \textit{train} selesai, \textit{snapshot} akan diperbarui kemudian digunakan
% untuk membuat sebuah berkas baru berupa JSON yang menyimpan \textit{trained model} tersebut
% sehingga \textit{supertagger} yang dibangun dapat memuat model tersebut tanpa perlu melakukan
% \textit{train} kembali.

% \begin{figure}\centering
% 	\bgroup
% 	\catcode`!=\active \def!{\upshape}
% 	\tabbedShortstack[l]{
% 		!{$\{$} &\\
% 			& !{\say{starting\_point}: $70$,}\\
% 			& !{\say{rows}:
% 				\{\say{$token_1$}: $1$, \say{$token_2$}: $2$, $\dots$, \say{$token_n$}: $N$\},}\\
% 			& !{\say{cols}:
% 				 \{\say{$token_1$}: $1$, \say{$token_2$}: $2$, $\dots$, \say{$token_m$}: $M$\},}\\
% 			& !{\say{values}: [}\\
% 			& !{\ \ \ \ [$value_{1,1}$, $value_{1,2}$, $\dots$, $value_{1,m}$],}\\
% 			& !{\ \ \ \ [$value_{2,1}$, $value_{2,2}$, $\dots$, $value_{2,m}$],}\\
% 			& \\
% 			& \\
% 			& \ \ \ \ $\vdots$\\
% 			& \\
% 			& \\
% 			& !{\ \ \ \ [$value_{n,1}$, $value_{n,2}$, $\dots$, $value_{n,m}$]}\\
% 			& !{]}\\
% 		!{$\}$} &
% 	}
% 	\egroup
% 	\caption{Format JSON untuk menyimpan \textit{snapshot}.}
% 	\label{snapshot:format}
% \end{figure}

% Pada Gambar \ref{snapshot:format}, \textit{field} \say{$starting\_point$} menyimpan informasi berupa
% bilangan cacah untuk memberikan tanda pada elemen ke-$i$ sebaiknya proses \textit{train} dimulai.
% Dalam hal ini, berdasarkan contoh, sebaiknya dimulai dari elemen ke-$70$.
% Selanjutnya, \textit{field} \say{$rows$} dan \say{$cols$} masing-masing menyatakan daftar baris dan
% kolom beserta nomor \textit{index}-nya untuk digunakan oleh \textit{field} \say{$values$} nantinya.
% \textit{Field} \say{$values$} merupakan representasi dari tabel yang memiliki $N$ baris serta $M$
% kolom.
% Adapun \say{$value$} pada elemen ke-$(i, j)$ menyimpan informasi berupa objek untuk model MaxEnt
% setelah proses \textit{train} pada \say{$starting\_point$} saat itu.
% Objek yang dimaksud berupa \say{konteks} yaitu berisikan daftar $k$ token sebelum dan sesudah token
% ke-$(i, j)$ juga daftar $l$ \textit{supertag} sebelum dan sesudah \textit{supertag} ke-$(i, j)$.
% Bahkan, \say{$value$} dapat pula meyimpan \textit{prefix} dan \textit{suffix} dari token ke-
% $(i, j)$ tersebut.

% \subsection{Membangun Library Haskell}

% Implementasi \textit{supertagger} bahasa Indonesia ini akan ditulis dalam bahasa pemrograman Haskell.
% Haskell dipilih karena kemampuannya dalam mengolah data teks terutama kegiatan \textit{parsing} berkat
% desain bahasa serta dukungan \textit{library} yang telah tersedia.
% Secara desain bahasa, Haskell sangat dekat dengan \textit{category theory} sehingga diharapkan
% \textit{supertag} yang disimpan di dalam program nantinya dapat berupa \textit{algebraic data type}
% yang semirip mungkin dengan \textit{category theory}.
% Hal ini menjadi pertimbangan karena tujuan akhir dari riset-riset yang berhubungan dengan CCG adalah
% untuk membangun CCG \textit{parser}.

% Pembangunan \textit{library} ini sangat penting karena nantinya \textit{library} ini dapat dimanfaatkan
% untuk membangun program \textit{command line interface} (CLI), \textit{web application}, dan sebagainya.
% Dalam hal ini, \textit{supertagger} yang dibangun dalam tugas akhir ini akan memiliki dukungan RESTful
% API agar tidak ada batasan dalam bahasa pemrograman yang harus digunakan untuk memanfaatkan
% \textit{supertagger} ini.
% Sangat mungkin pula nantinya terdapat program yang memanfaatkan \textit{library} Haskell ini secara
% langsung di programnya.
% Demikian itu, menyediakan \textit{library} untuk suatu bahasa pemrograman (dalam hal ini, Haskell)
% dirasa perlu.

% \subsection{Menguji Model MaxEnt}

% Pengujian model MaxEnt setelah proses \textit{train} dilakukan dengan cara mencocokkan \textit{supertag}
% yang dihasilkan oleh \textit{supertagger} dengan \textit{supertag} yang telah diberikan secara manual
% terhadap suatu kalimat yang diberikan.
% Dengan demikian, setidaknya akan ada dua \textit{dataset} yaitu \textit{data train} dan
% \textit{data test}.
% \textit{Data train} dibuat secara semi-otomatis sesuai dengan bagian \say{Pembuatan Dataset}.
% Selanjutnya, \textit{data test} dibuat secara manual yang mana berjumlah $20\%$ hingga $30\%$ dari
% jumlah kalimat pada \textit{data train}.
% Sebagai catatan, untuk kata atau token yang tidak ada di \textit{data train} tetapi ada di
% \textit{data test} (\textit{unseen words}) akan diberikan \textit{supertag} berupa
% \textit{category} $\text{N}$.

% \section{Membangun CCG Supertagger Versi CLI}

% Sebagai langkah awal, program CLI akan dibangun lebih dahulu sebelum versi \textit{web application}-nya
% rilis.
% Hal ini agar validasi dapat dilakukan yaitu untuk memeriksa apakah keluaran dari \textit{supertagger}
% sudah sesuai.
% Adapun \textit{flowchart} dari program CLI \textit{supertagger} ini terdapat pada Gambar
% \ref{flowchart:supertagger:1}.

% \begin{figure}[h!]\centering
% 	\begin{tikzpicture}[node distance=2cm]
% 		\node (start) [cloud] {Start};
% 		\node (in) [io, below of=start, yshift=0.5cm] {\textit{Input} teks bahasa Indonesia};
% 		\node (pro1) [process, below of=in, yshift=0.5cm] {\textit{Tokenize}};
% 		\node (pro2) [process, below of=pro1, yshift=0.5cm] {Jadikan \textit{queue}};
% 		\node (dec) [decision, below of=pro2, yshift=-1cm] {\textit{Queue} kosong?};
% 		\node (pro3) [process, below of=dec, yshift=-1cm] {Menghitung probabilitas terhadap token};
% 		\node (pro4) [process, below of=pro3, yshift=0.5cm] {Memilih \textit{supertag} yang sesuai};
% 		\node (out) [io, below of=pro4, yshift=0.5cm] {\textit{Output} CCG \textit{supertag}};
% 		\node (stop) [cloud, below of=out, yshift=0.5cm] {Stop};
% 		\draw [arrow] (start) -- (in);
% 		\draw [arrow] (in) -- (pro1);
% 		\draw [arrow] (pro1) -- (pro2);
% 		\draw [arrow] (pro2) -- (dec);
% 		\draw [arrow] (dec) -- node[anchor=east]{Tidak} (pro3);
% 		\draw [arrow] (pro3) -- (pro4);
% 		\draw [arrow] (pro4) -- +(-4,0) |- (dec);
% 		\draw [arrow] (dec) -- node[anchor=south]{Ya} +(4,0) |- (out);
% 		\draw [arrow] (out) -- (stop);
% 	\end{tikzpicture}
% 	\caption{Alur kerja proses \textit{tagging}.}
% 	\label{flowchart:supertagger:1}
% \end{figure}

% \subsection{Tokenize}

% Pertama-tama, \textit{toknize} akan memisahkan antar kalimat yang dipisahkan oleh karakter titik
% (\say{.} tanpa tanda kutip) yang diikuti oleh spasi. Selanjutnya, untuk setiap kalimat,
% \textit{tokenize} akan memisahkan beberapa bagian yang terindikasi sebagai token.
% Token umumnya berupa kata yang dipisahkan oleh spasi.
% Namun, token juga dapat berupa karakter khusus seperti karakter tanda koma (\say{,} tanpa tanda kutip).
% Dengan demikian, \textit{supertagger} dapat dengan mudah melakukan kalkulasi probabilitas untuk
% masing-masing token.
% Proses \textit{tokenize} pada awalnya terlihat mudah.
% Sayangnya, proses ini cukup \textit{tricky} karena ada beberapa hal tidak umum yang seringkali muncul.
% Sebagai contoh, \say{Rp5.000,00} memiliki dua token yaitu \say{Rp} sebagai \textit{currency} dan
% \say{5.000,00} sebagai nominalnya.

% \subsection{Menghitung Probabilitas}

% Probabilitas untuk sebuah token ke-$(i, j)$ dapat dihitung dengan menggunakan \textit{beam search}.
% Persamaan yang digunakan adalah persamaan \ref{maxent:equation:4}.
% Persamaan tersebut memiliki sebuah \textit{category sequence} $C$ serta kalimat $S$ sebagai syarat
% dari probabilitas tersebut.

% \begin{equation}\label{maxent:equation:4}
% 	p(C|S) = \prod_i p(c_i, h_i)
% \end{equation}

% \noindent
% dimana $c_i$ adalah \textit{category} ke-$i$ dalam \textit{sequence} tersebut dan $h_i$ berupa konteks
% dari token ke-$i$.
% \textit{Beam search} digunakan untuk mendapatkan $N$ \textit{sequence} teratas saat proses
% \textit{tagging} dilakukan.
% Dalamn hal ini, $N = 10$ digunakan.

% \subsection{Memilih Supertag}

% Setelah proses penghitungan probabilitas selesai dilakukan, proses pemilihan \textit{supertag} dilakukan
% berdasarkan $N$ \textit{sequence} teratas yang disimpan.
% Berhubung MaxEnt yang digunakan adalah \textit{conditional maximum entorpy}, maka selanjutnya akan
% dilakukan proses pemilihan \textit{supertag} berdasarkan kondisi-kondisi yang terpenuhi.

CCGtown dibangun dengan menggunakan bahasa pemrograman Python dan JavaScript.
Adapun \textit{framework} yang digunakan adalah Django.
Versi awal CCGtown merupakan sebuah \textit{proof-of-concept} dari
\textit{open source graphical annotation tool} berbasis web yang dilengkapi dengan fitur
penganotasian semi-otomatis.
Bahasa pemrograman Python digunakan karena sebagian besar \textit{library} untuk CCG
sudah tersedia di PyPi \footnote{\url{https://pypi.org/}}.
Salah satu \textit{library} penting yang digunakan sebagai dasar dari fitur penganotasian
semi-otomatis adalah NTLK \footnote{\url{http://www.nltk.org/}}.
Selanjutnya, Django digunakan untuk mempercepat proses pengembangan aplikasi.
Adapun JavaScript digunakan untuk menjadikan CCGtown aplikasi berbasis web yang interaktif.

Alur kerja CCGtown pada umumnya adalah (1) pengguna melakukan registrasi, (2) pengguna
melakukan \textit{login} ke sistem, (3) pengguna membuat proyek baru, (4) pengguna
menambahkan kalimat yang ingin dianotasi, (5) pengguna melakukan anotasi kemudian melakukan
\textit{generate} CCG \textit{derivation} dan/atau melakukan modifikasi
\textit{derivation}-nya apabila diperlukan, dan (6) pengguna melakukan \textit{export}
setelah selesai melakukan anotasi. Alur kerja tersebut mempengaruhi desain sistem dari
CCGtown. Salah satunya adalah desain dari \textit{database} yang akan digunakan.


\section{Desain Database}
CCGtown menggunakan PostgreSQL sebagai
DBMS\footnote{Database Management System}-nya.
Hal ini karena PostgreSQL memiliki kemampuan untuk menyimpan struktur data
JSON\footnote{JavaScript Object Notation} sehingga memudahkan CCGtown untuk menyimpan
format JSON dari CCG \textit{derivation} yang telah dimanipulasi oleh pengguna melalui
fitur \textit{editable CCG derivation}.
PostgreSQL juga memiliki banyak fitur lain termasuk di antaranya dukungan
dari \textit{non-relational database model} (seperti \textit{multi-model graph})
sehingga apabila di waktu yang akan datang CCGtown memerlukan perubahan signifkan
terhadap desain \textit{database}-nya tidak perlu mengganti DBMS yang digunakan.
Fitur lain seperti \textit{function} dan \textit{procedure} juga akan sangat membantu
pengembangan CCGtown di waktu yang akan datang.

CCGtown versi awal sejatinya hanya membutuhkan tiga tabel saja yaitu tabel
\textit{accounts} untuk menyimpan pengguna yang terdaftar, tabel
\textit{projects} untuk menyimpan proyek-proyek yang sudah dibuat, dan tabel
\textit{sentences} untuk menyimpan kalimat-kalimat yang akan dianotasikan.
Tiga tabel tersebut sudah cukup untuk membangun \textit{proof-of-concept} dari
alat anotasi CCG yang akan dibangun. Adapun
ERD\footnote{Entity Relationship Diagram}-nya dapat dilihat pada Gambar\ref{erd-1}.

\begin{figure}\centering\small
  \scalebox{.75}{
  \begin{tikzpicture}[node distance=1.5cm, every edge/.style={link}]
    \node[entity] (acc) {Accounts};
    \node[attribute] (acc-id) [above=of acc] {\key{ID}} edge (acc);
    \node[attribute] (acc-uuid) [above right=of acc] {\key{UUID}} edge (acc);
    \node[attribute] (acc-email) [right=of acc] {Email} edge (acc);
    \node[attribute] (acc-password) [below right=of acc] {Password} edge (acc);

    \node[relationship] (creates) [left=of acc] {Creates} edge (acc);

    \node[entity] (prj) [below=of creates] {Projects} edge (creates);
    \node[attribute] (prj-id) [above right=of prj] {\key{ID}} edge (prj);
    \node[attribute] (prj-uuid) [above left=of prj] {\key{UUID}} edge (prj);
    \node[attribute] (prj-author) [left=of prj] {Author ID} edge (prj);
    \node[attribute] (prj-name) [below left=of prj] {Name} edge (prj);
    \node[attribute] (prj-status) [below=of prj] {Status} edge (prj);
    \node[attribute] (prj-rules) [below right=of prj] {Rules} edge (prj);

    \node[relationship] (adds) [right=0.5cm and 2cm of prj] {Adds} edge (prj);

    \node[entity] (snt) [below=of adds] {Sentences} edge (adds);
    \node[attribute] (snt-id) [left=of snt] {\key{ID}} edge (snt);
    \node[attribute] (snt-uuid) [below left=of snt] {\key{UUID}} edge (snt);
    \node[attribute] (snt-project) [below=of snt] {Project ID} edge (snt);
    \node[attribute] (snt-words) [below right=of snt] {Words} edge (snt);
    \node[attribute] (snt-cats) [right=of snt] {Categories} edge (snt);
    \node[attribute] (snt-deriv) [above right=of snt] {Derivations} edge (snt);
  \end{tikzpicture}
  }
  \caption{Conceptual Entity Relationship Diagram (ERD) CCGtown}
  \label{erd-1}
\end{figure}

Masing-masing tabel memiliki dua \textit{key} yaitu $ID$ dan
$UUID$\footnote{Universally Unique IDentifier}.
$ID$ merupakan \textit{primary key} \textit{integer} dengan \textit{auto increment}
yang berfungsi sebagai \textit{identifier} untuk melakukan operasi
\textit{update} maupun \textit{delete}.
Adapun $UUID$ merupakan \textit{indexed column} yang berfungsi sebagai
\textit{indentifier} publik (dapat dilihat oleh pengguna melalui URL)
yang mana digunakan untuk operasi \textit{read}.
$ID$ tidak digunakan sebagai \textit{identifier} publik karena pengguna dapat
melakukan \textit{brute-force} untuk mencari proyek ataupun kalimat berdasarkan
$ID$ yang bukan miliknya.
Demikian itu alasan ditambahkannya atribut $UUID$.
Alasan kenapa CCGtown tetap menyimpan kolom $ID$ adalah karena $ID$ nantinya akan
digunakan untuk membuat \textit{pagination}.

Pada tabel $accounts$, selain $ID$ dan $UUID$ juga memiliki atribut $email$ dan
$password$. Masing-masing atribut tersebut menggunakan tipe data \textit{string}
atau VARCHAR di PostgreSQL.
Tabel $accounts$ memiliki hubungan \textit{one-to-many} terhadap tabel $projects$.
Adapun atribut tabel $projects$ adalah $author\_id$, $name$, $status$, dan $rules$.
Atribut $author\_id$ merupakan \textit{foreign key} (\textit{indexed}) yang
mengarah kepada tabel $accounts$ dan tipe data yang digunakan sama dengan
atribut $ID$ yang terdapat di tabel $accounts$.
Atribut $name$ menggunakan tipe data \textit{string} (VARCHAR).
Atribut $status$ menggunakan tipe data \textit{integer} yang berperan sebagai
\textit{enum} ($0$ = \textit{just created}, $1$ = \textit{in progress},
$2$ = \textit{finished}, dan $3$ = \textit{dropped}).
Tabel $projects$ memiliki hubungan \textit{one-to-many} terhadap tabel $sentences$.
Adapun atribut tabel $sentences$ adalah $project\_id$, $words$, $categories$, dan
$derivations$. Atribut $project\_id$ merupakan \textit{foreign key}
(\textit{indexed}) yang mengarah kepada tabel $projects$ dan tipe data yang digunakan
sama dengan atribut $ID$ yang terdapat di tabel $projects$.
Sisanya, atribut $words$, $categories$, dan $derivations$ menggunakan tipe data JSON.


\section{Desain Sistem}
CCGtown sejatinya memiliki desain sistem yang cukup sederhana.
Fungsionalitas yang akan didukung untuk versi awal adalah (1) \textit{register} dan
\textit{login}, (2) manajemen proyek (CRUD\footnote{Create, Read, Update, Delete}
), (3) dan manajemen kalimat (CRUD).
Pada manajemen kalimat, CCGtown menggunakan JavaScript untuk membuat pembuatan
maupun perubahan CCG \textit{derivation} menjadi lebih interaktif.
Selain tiga fungsionalitas tersebut, CCGtown juga menambahkan fungsionalitas tambahan
seperti \textit{auto-assign category} yang dilakukan di sisi \textit{frontend}.
Kemudian, CCGtown juga menambahkan fungsionalitas tambahan di sisi \textit{backend}
yaitu CCG \textit{derivation generator} dengan memanfaatkan \textit{library} NLTK
dan kemampuan untuk melakukan \textit{export} CCG \textit{derivation} yang disimpan
di \textit{database}.

Pengguna harus terdaftar terlebih dahulu sebelum dapat melaukan anotasi sehingga
langkah awal yang harus dibangun adalah fungsionalitas \textit{register}.
Alur proses pendaftaran pengguna dapat dilihat pada Gambar \ref{flowchart:register}.
Berhubung fokus saat ini adalah \textit{proof-of-concept}, informasi yang dibutuhkan
untuk mendaftar hanyalah \textit{email} dan \textit{password}. Adapun
\textit{password confirmation} digunakan untuk memvalidasi \textit{password}
sehingga dapat mengurangi risiko pengguna melupakan
\textit{password}-nya yang baru saja di-\textit{input}.
Saat pengguna melakukan pendaftaran, sistem akan memeriksa apakah \textit{email} yang
didaftar sudah terdapat di \textit{database}.
Apabila sudah terdaftar, pengguna akan dialihkan ke halaman \textit{register} kembali
dan mendapatkan \textit{flash message} dengan keterangan "email sudah terdaftar".
Sebaliknya, sistem akan melakukan \textit{input} data tersebut ke dalam \textit{database}
lalu mengalihkan pengguna ke halaman \textit{login}.
Ketika dialihkan ke halaman \textit{login}, pengguna akan melihat \textit{flash message}
dengan keterangan "pengguna berhasil didaftarkan".
Pada tahap ini pengguna sudah dapat melakukan \textit{login} ke dalam sistem CCGtown.

\begin{figure}\centering\small
  \scalebox{.75}{
	\begin{tikzpicture}[node distance=2cm]
    \node (start) [cloud] {Start};
    \node (input) [io, below of=start, yshift=0.5cm] {\textit{Input email, password,} dan \textit{password confirmation}};
    \node (check) [process, below of=input, yshift=0.5cm] {Mencari pengguna berdasarkan \textit{email}};
    \node (is-registered) [decision, below of=check, yshift=-1.25cm] {\textit{Email} sudah terdaftar?};
    \node (registered) [process, right of=is-registered, yshift=0cm, xshift=4.5cm] {\textit{Redirect} ke halaman \textit{register}};
    \node (registering) [process, below of=is-registered, yshift=-1.75cm] {\textit{Input} informasi pengguna ke \textit{database}};
    \node (redirect) [process, below of=registering, yshift=0.5cm] {\textit{Redirect} ke halaman \textit{login}};
    \node (stop) [cloud, below of=redirect, yshift=0.5cm] {Stop};

    \draw [arrow] (start) -- (input);
    \draw [arrow] (input) -- (check);
    \draw [arrow] (check) -- (is-registered);
    \draw [arrow] (is-registered) -- node[anchor=south]{Ya} (registered);
    \draw [arrow] (registered) -- +(3,0) |- (input);
    \draw [arrow] (is-registered) -- node[anchor=east]{Tidak} (registering);
    \draw [arrow] (registering) -- (redirect);
    \draw [arrow] (redirect) -- (stop);
  \end{tikzpicture}
  }
	\caption{Alur proses pendaftaran pengguna.}
  \label{flowchart:register}
\end{figure}

Pada proses "\textit{input} informasi pengguna ke \textit{database}" CCGtown melakukan
\textit{password hashing} dengan menggunakan Bcrypt.
Informasi sensitif seperti \textit{password} sebaiknya tidak disimpan sebagai
\textit{plain text}. Demikian itu CCGtown menggunakan \textit{password hashing}.
Apabila hal buruk terjadi seperti misalnya \textit{data breach} (kebocoran data),
\textit{password} pengguna tidak dapat langsung digunakan.
Peretas perlu mencari cara untuk memecahkan \textit{password} tersebut.
Bcrypt merupakan skema \textit{password hashing} berbasis Blowfish \textit{block cipher}
yang didesain untuk lebih \textit{resistant} terhadap serangan \textit{brute-force}
\cite{bcrypt}.
Serangan \textit{brute-force} merupakan upaya peretas untuk menebak \textit{password}
dengan cara membuat \textit{wordlist} yang kemudian dicocokkan dengan \textit{hash}
yang terbentuk satu-demi-satu.
Meskipun terjadi \textit{data breach}, peretas perlu usaha ekstra untuk dapat menebak
\textit{password} dari satu pengguna.
Hal ini mengurangi kerugian yang akan dialami oleh CCGtown apabila \textit{data breach}
benar-benar terjadi.

Selanjutnya, setelah melakukan registrasi, pengguna dapat melakukan \textit{login} ke
sistem CCGtown.
Proses yang dilakukan pada umumnya sama dengan aplikasi web yang memiliki
kemampuan \textit{register} dan \textit{login}. Alur proses \textit{login} dapat dilihat
pada Gambar \ref{flowchart:login}.
Setelah pengguna melakukan \textit{input} \textit{email} dan \textit{password}-nya,
CCGtown akan melakukan pencarian di \textit{database} apakah \textit{email} yang
diberikan terdaftar. Apabila tidak terdaftar, pengguna akan dialihkan ke halaman
\textit{login} dan diberikan \textit{flash message} "\textit{Email} dan/atau
\textit{password} tidak cocok". Pesan ini diberikan agar peretas tidak dapat mencari
tahu \textit{email} mana saja yang sudah terdaftar. Selanjutnya, apabila akun
dengan \textit{email} tersebut ada, maka langkah selanjutnya adalah mencocokkan
\textit{password} yang diberikan oleh pengguna dan \textit{password} yang telah
disimpan di \textit{database}. Kemudian, sistem melakukan Bcrypt \textit{sync}.
Apabila tidak berhasil, pengguna akan dialihkan ke halaman \textit{login}
dan diberikan \textit{flash message} "\textit{Email} dan/atau \textit{password}
tidak cocok". Sebaliknya, pengguna akan dialihkan ke halaman Projects yang berisi
daftar proyek yang telah dibuat sebelumnya.

\begin{figure}\centering\small
  \scalebox{.75}{
	\begin{tikzpicture}[node distance=2cm]
    \node (start) [cloud] {Start};
    \node (input) [io, below of=start, yshift=0.5cm] {\textit{Input email} dan \textit{password}};
    \node (check) [process, below of=input, yshift=0.5cm] {Mencari pengguna berdasarkan \textit{email}};
    \node (is-registered) [decision, below of=check, yshift=-1.25cm] {\textit{Email} sudah terdaftar?};
    \node (not-registered) [process, left of=is-registered, yshift=0cm, xshift=-4.5cm] {\textit{Redirect} ke halaman \textit{login}};
    \node (logging-in) [process, below of=is-registered, yshift=-1.75cm] {Mencocokkan \textit{password} dengan Bcrypt \textit{sync}};
    \node (is-matched) [decision, below of=logging-in, yshift=-1.25cm] {\textit{Password} cocok?};
    \node (redirect) [process, below of=is-matched, yshift=-1.75cm] {\textit{Redirect} ke halaman Projects};
    \node (stop) [cloud, below of=redirect, yshift=0.5cm] {Stop};

    \draw [arrow] (start) -- (input);
    \draw [arrow] (input) -- (check);
    \draw [arrow] (check) -- (is-registered);
    \draw [arrow] (is-registered) -- node[anchor=south]{Tidak} (not-registered);
    \draw [arrow] (not-registered.north) -- +(0,0) |- (input);
    \draw [arrow] (is-registered) -- node[anchor=east]{Ya} (logging-in);
    \draw [arrow] (logging-in) -- (is-matched);
    \draw [arrow] (is-matched.west) -| +(-4.6,0) -- node[anchor=east]{Tidak} (not-registered.south);
    \draw [arrow] (is-matched) -- node[anchor=east]{Ya} (redirect);
    \draw [arrow] (redirect) -- (stop);
  \end{tikzpicture}
  }
	\caption{Alur proses \textit{login} ke sistem CCG.}
  \label{flowchart:login}
\end{figure}


\section{Pertimbangan UI/UX}
TBA.

%

%%%%%%%%%%%%%%%%%%%%%%%%%%%%%%%%%%%%%%%%%%%%%%%%%%%%%%%%%%%%%%%
% Berikut adalah untuk membuat Daftar pustaka
% Style daftar pustaka dapat disesuaikan dengan mengubah 
%\bibliographystyle{kode}
% dengan kode = acm maka hasinya contoh [1]
% dengan kode = agsm maka hasilnya Harvard style (Sukimin, 2017)
%%%%%%%%%%%%%%%%%%%%%%%%%%%%%%%%%%%%%%%%%%%%%%%%%%%%%%%%%%%%%%%%
\cleardoublepage
\addcontentsline{toc}{chapter}{Daftar Pustaka}
\bibliographystyle{plain} %harvard style
\bibliography{References}
%
%\pagebreak
%%%%%%%%%%%%%%%%%%%%%%%%%%%%%%%%%%%%%%%%%%%%%%%%%%%%%%%%%%%%%%%
% Berikut untuk lampiran
%%%%%%%%%%%%%%%%%%%%%%%%%%%%%%%%%%%%%%%%%%%%%%%%%%%%%%%%%%%%%%
\cleardoublepage
\addcontentsline{toc}{chapter}{Lampiran}
\chapter*{Lampiran}

\end{document}
